% =============================================================
%  Tugas Individu 09 — Menyejajarkan Ruang Embedding
%  IF25-40304 · Pembelajaran Mesin Multimodal · ITERA
% =============================================================
\documentclass[12pt, a4paper]{article}

\usepackage{tugas-style}

\renewcommand{\assignmentnumber}{09}
\renewcommand{\doctitle}{Tugas Individu 09}
\renewcommand{\duedate}{\color{ctpBlue}\faCalendarCheck\ \ Tenggat: sesuai jadwal kelas \faCalendarCheck}

\begin{document}

% ── HEADER ──────────────────────────────────────────────────
\makeassignmentheader

\hrule width0.3\textwidth
\revisioninfo{%
  \vhEntry{0.1}{23 September 2026}{I.W.W.}{Draf inisial tugas.}
}
\hrule
\vskip .3in

\tableofcontents
\clearpage

% =============================================================
\section{Tujuan Pembelajaran}\label{sec:objectives}

Setelah menyelesaikan tugas ini, mahasiswa mampu:
\begin{objectives}
  \item Mengukur keselarasan dua modalitas secara kuantitatif melalui matriks kesamaan dan nilai \emph{InfoNCE}.
  \item Menyelaraskan dua ruang \emph{embedding} dan membuktikan perbaikan kinerjanya dengan angka.
  \item Menjelaskan pengaruh suhu (\code{tau}) terhadap tujuan kontrastif dan kestabilan pelatihan.
  \item Mengukur kinerja pencarian lintas modal dengan metrik \emph{Recall@K}.
  \item Menilai secara kritis batas arsitektur \emph{dual-encoder} dan mengujinya pada kasus relasi komposisional.
\end{objectives}

% =============================================================
\section{Catatan dan Batasan}\label{sec:notes}
\vspace{-.1cm}

\begin{minipage}{\linewidth}
\begin{bclogo}[couleur=ctpBase, arrondi=0.1, logo=\bccrayon, ombre=false, couleurBord=ctpSurface, epBord=0.6]{Perhatikan hal berikut:}
  \begin{coloredPen}
    \item Tugas ini bersifat \textbf{individu}; tidak diperkenankan bekerja sama atau menyalin jawaban.
    \item Soal 1--4 memakai data \textbf{sintetis} dengan \emph{seed} tetap, sehingga hasil dapat direproduksi tanpa mengunduh apa pun.
    \item Soal 5b memakai model \textbf{CLIP pra-latih}. Bobot model \textbf{tidak boleh} diubah --- seluruh Soal 5b bersifat inferensi.
    \item Laporkan selalu \textbf{baseline sebelum penyelarasan}. Angka setelah penyelarasan tidak bermakna tanpa pembanding.
    \item Gunakan \code{numpy} untuk Soal 1--4 dan \code{torch} + \code{transformers} untuk Soal 5b.
  \end{coloredPen}
\end{bclogo}
\end{minipage}

\begin{minipage}{\linewidth}
\begin{bclogo}[couleur=ctpMantle, arrondi=0.1, logo=\bcinfo, ombre=false, couleurBord=ctpSurface, epBord=0.6]{Ekspektasi Hasil}
  Dengan \code{SEED = 42} dan data sintetis pada tugas ini, hasil yang wajar adalah:

  \begin{itemize}[itemsep=1pt,parsep=0pt,topsep=2pt,partopsep=0pt]
    \item Sebelum penyelarasan: \emph{InfoNCE} berada di kisaran \textbf{10--11} dan \emph{Recall@1} sekitar \textbf{1/6}.
    \item Sesudah penyelarasan: \emph{InfoNCE} turun di bawah \textbf{0,2} dan \emph{Recall@1} mencapai \textbf{6/6}.
    \item Suhu mengecil pada data yang sudah selaras menurunkan \emph{loss}, tetapi suhu yang terlalu kecil membuat pelatihan tidak stabil.
  \end{itemize}

  Bila hasil Anda menyimpang jauh (misalnya \emph{Recall@1} tetap 1/6 setelah penyelarasan), periksa: data belum dipusatkan sebelum dekomposisi nilai singular; sumbu matriks kesamaan tertukar; atau vektor belum dinormalisasi sebelum dihitung kesamaan kosinus.
\end{bclogo}
\end{minipage}

% =============================================================
\section{Perangkat Lunak dan Peralatan yang Diperlukan}\label{sec:requiredsw}

\begin{itemize}[itemsep=2pt,parsep=0pt,topsep=2pt,partopsep=2pt]
  \item[\color{ctpBlue}\faPython] \textbf{Python} 3.8 atau lebih baru.
  \item[\color{ctpBlue}\faFileCode] \textbf{Paket wajib:} \code{numpy} (Soal 1--4), \code{torch} + \code{transformers} + \code{pillow} (Soal 5b).
  \item[\color{ctpBlue}\faDownload] \textbf{Data:} sintetis untuk Soal 1--4; satu gambar pilihan sendiri untuk Soal 5b.
  \item[\color{ctpBlue}\faFileCode] \textbf{Perangkat:} CPU memadai untuk Soal 1--4; Soal 5b disarankan memakai GPU bila tersedia (bukan syarat).
\end{itemize}

\begin{codewindow}{setup.sh}
\begin{lstlisting}[style=pycode]
pip install numpy torch transformers pillow
\end{lstlisting}
\end{codewindow}

\begin{minipage}{\linewidth}
\begin{bclogo}[couleur=ctpMantle, arrondi=0.1, logo=\bcinfo, ombre=false, couleurBord=ctpSurface, epBord=0.6]{Batasan Cakupan}
  Soal 1--4 memakai data multimodal \textbf{sintetis} yang dibangkitkan secara programatik.
  Kedua modalitas memandang faktor laten yang sama, tetapi modalitas kedua melihatnya
  melalui rotasi ortogonal. Konstruksi ini membuat penyelarasan dapat dihitung dalam
  bentuk tertutup, sehingga Anda dapat memisahkan persoalan \emph{alignment} dari
  persoalan \emph{pelatihan jaringan}.
\end{bclogo}
\end{minipage}

% =============================================================
\section{Pernyataan Masalah}\label{sec:intro}

Sebuah katalog berisi ratusan ribu foto produk. Pelanggan mengetik deskripsi bebas, dan
sistem harus mengembalikan foto yang paling sesuai. Tidak tersedia data berpasangan
berlabel, dan latensi pencarian harus di bawah satu detik.

Ini adalah persoalan \textbf{penyejajaran lintas modal}: bagaimana membuat satu gambar
dan satu kalimat yang membicarakan hal yang sama berada pada titik yang berdekatan di
ruang vektor. Sesi Pertemuan 09 menunjukkan bahwa model \emph{dual-encoder} seperti CLIP
melakukannya tanpa pernah mempertemukan lapisan kedua modalitas.

Anda akan membuktikan mekanismenya dengan angka: mula-mula tanpa penyelarasan, kemudian
dengan penyelarasan, lalu pada model pra-latih.

% =============================================================
\section{Persyaratan}\label{sec:requirements}

Perhatikan setiap persyaratan pada soal-soal berikut.

% ─────────────────────────────────────────────────────────────
\subsection{Soal 1 --- Pasangan Sintetis dan Baseline InfoNCE}

Sebelum menyelaraskan apa pun, kita harus mengukur seberapa buruk keadaan awalnya.
Inilah \emph{baseline} yang menjadi pembanding sepanjang tugas.

\begin{enumerate}
  \item Bangkitkan faktor laten bersama $\mathbf{z}$, lalu bentuk dua modalitas yang
        memandang faktor tersebut melalui proyeksi berbeda. Pusatkan $\mathbf{z}$
        terhadap rata-ratanya.
  \item Hitung matriks kesamaan kosinus antara kedua modalitas, lalu tampilkan nilai
        diagonalnya (pasangan benar) dan rerata di luar diagonal.
  \item Hitung nilai \emph{InfoNCE} dan \emph{Recall@1} pada keadaan ini.
  \item Jawab dalam 3--5 kalimat: mengapa matriks kesamaan tampak acak sebelum
        penyelarasan, meskipun kedua modalitas sebenarnya mengandung informasi yang sama?
\end{enumerate}

\begin{codewindow}{soal1\_data.py}
\begin{lstlisting}[style=pycode]
import numpy as np

SEED, N, DIM, TAU = 42, 6, 4, 0.10
rng = np.random.default_rng(SEED)

# Faktor laten bersama, dipusatkan
z = rng.normal(0.0, 1.0, size=(N, DIM))
z = z - z.mean(axis=0, keepdims=True)

# Modalitas kedua melihat faktor yang sama lewat rotasi ortogonal Q
Q, _ = np.linalg.qr(rng.normal(0.0, 1.0, size=(DIM, DIM)))

gambar = z + 0.10 * rng.normal(size=(N, DIM))
teks = z @ Q.T + 0.10 * rng.normal(size=(N, DIM))

def normalisasi(X):
    return X / np.linalg.norm(X, axis=1, keepdims=True)

def matriks_kesamaan(A, B):
    return normalisasi(A) @ normalisasi(B).T

def infonce(S, tau=TAU):
    Z = S / tau
    Z = Z - Z.max(axis=1, keepdims=True)          # stabilkan secara numerik
    log_softmax = Z - np.log(np.exp(Z).sum(axis=1, keepdims=True))
    return -np.diag(log_softmax).mean()

def recall_at_1(S):
    return int(np.mean(np.argmax(S, axis=1) == np.arange(len(S))) * len(S))

S0 = matriks_kesamaan(gambar, teks)
print("diagonal :", np.diag(S0))
print("rerata luar diagonal: %+.3f" % S0[~np.eye(N, dtype=bool)].mean())
print("InfoNCE  : %.4f" % infonce(S0))
print("Recall@1 : %d/%d" % (recall_at_1(S0), N))
\end{lstlisting}
\end{codewindow}

% ─────────────────────────────────────────────────────────────
\subsection{Soal 2 --- Penyelarasan Bentuk Tertutup}

Sekarang kita selaraskan ruang teks terhadap ruang gambar. Karena hubungan antar
keduanya ortogonal, penyelarasan optimal dapat dihitung langsung lewat dekomposisi
nilai singular.

\begin{enumerate}
  \item Hitung matriks rotasi optimal $\mathbf{R}$, lalu selaraskan vektor teks
        dengan $\mathbf{R}$ tersebut.
  \item Hitung ulang matriks kesamaan, \emph{InfoNCE}, dan \emph{Recall@1}.
  \item Sajikan perbandingan sebelum dan sesudah dalam satu tabel.
  \item Jawab dalam 4--6 kalimat: apa yang sesungguhnya berubah pada langkah ini, dan
        mengapa hal itu cukup untuk memperbaiki kinerja pencarian?
\end{enumerate}

\begin{codewindow}{soal2\_procrustes.py}
\begin{lstlisting}[style=pycode]
M = teks.T @ gambar
U, s, Vt = np.linalg.svd(M)
R = U @ Vt
teks_selaras = teks @ R

S1 = matriks_kesamaan(gambar, teks_selaras)
print("diagonal :", np.diag(S1))
print("rerata luar diagonal: %+.3f" % S1[~np.eye(N, dtype=bool)].mean())
print("InfoNCE  : %.4f" % infonce(S1))
print("Recall@1 : %d/%d" % (recall_at_1(S1), N))
\end{lstlisting}
\end{codewindow}

% ─────────────────────────────────────────────────────────────
\subsection{Soal 3 --- Pengaruh Suhu}

Suhu menentukan seberapa tajam model membedakan kandidat.

\begin{enumerate}
  \item Hitung \emph{InfoNCE} pada data yang sudah selaras untuk setidaknya tiga
        nilai suhu, misalnya $0{,}05$, $0{,}10$, dan $0{,}50$.
  \item Tampilkan hasilnya dalam bentuk tabel.
  \item Jawab dalam 3--5 kalimat: mengapa suhu yang sangat kecil menurunkan
        \emph{loss}, tetapi tetap tidak dianjurkan dipakai pada pelatihan?
\end{enumerate}

\begin{codewindow}{soal3\_suhu.py}
\begin{lstlisting}[style=pycode]
for t in (0.05, 0.10, 0.50):
    print("tau=%.2f: InfoNCE=%.4f" % (t, infonce(S1, tau=t)))
\end{lstlisting}
\end{codewindow}

% ─────────────────────────────────────────────────────────────
\subsection{Soal 4 --- Retrieval dan Recall@K}

Pada tugas nyata, yang diukur bukan \emph{loss} melainkan peringkat hasil pencarian.

\begin{enumerate}
  \item Untuk setiap kueri teks, tentukan peringkat gambar yang benar.
  \item Laporkan \emph{Recall@1} dan jelaskan mengapa pada katalog besar angka ini
        dapat terlihat rendah meskipun modelnya baik.
  \item Sebutkan satu cara agar \emph{Recall@K} tinggi tanpa melatih ulang model.
\end{enumerate}

\begin{codewindow}{soal4\_retrieval.py}
\begin{lstlisting}[style=pycode]
for i in range(N):
    rank = np.argsort(-S1[i])
    pos = int(np.where(rank == i)[0][0]) + 1
    print("  teks %d: peringkat %d" % (i, pos))
\end{lstlisting}
\end{codewindow}

% ─────────────────────────────────────────────────────────────
\clearpage
\subsection{Soal 5 --- Analisis Kritis}

Jawab pertanyaan berikut dalam bentuk paragraf singkat (masing-masing 5--8 kalimat).

\textbf{(a) Perbandingan Pola Arsitektur.}

Buat tabel yang membandingkan tiga pola arsitektur multimodal (\emph{dual-stream},
\emph{multi-encoder}, dan \emph{dual-encoder}) pada tiga aspek: tempat modalitas
bertemu, kedalaman interaksi yang mungkin dipelajari, dan kesesuaiannya untuk
pencarian berskala besar. Pilih satu pola untuk pernyataan masalah pada
Bagian~\ref{sec:intro} dan pertanggungjawabkan pilihan Anda secara eksplisit.

\textbf{(b) Uji Kegagalan Komposisional (memerlukan kode).}

Uji apakah model pra-latih benar-benar menangkap \emph{relasi} antar objek atau hanya
kehadiran objek. Gunakan satu gambar pilihan Anda dan empat kandidat teks yang memakai
kata yang sama tetapi berbeda relasi atau negasi. Catat probabilitas setiap kandidat,
lalu jawab: apakah kegagalan yang Anda amati disebabkan oleh kurangnya data, kurangnya
kapasitas, atau tujuan pelatihannya? Jelaskan dasar penilaian Anda.

\begin{codewindow}{soal5\_komposisional.py}
\begin{lstlisting}[style=pycode]
import torch
from PIL import Image
from transformers import CLIPModel, CLIPProcessor

NAMA = "openai/clip-vit-base-patch32"
model = CLIPModel.from_pretrained(NAMA)
proc = CLIPProcessor.from_pretrained(NAMA)
model.eval()

gambar = Image.open("gambar_pilihan_anda.jpg").convert("RGB")
kandidat = [
    "sebuah sepeda di atas mobil",      # relasi benar
    "sebuah mobil di atas sepeda",      # relasi dibalik
    "sebuah sepeda dan sebuah mobil",   # tanpa relasi
    "sebuah mobil tanpa sepeda",        # negasi
]

inputs = proc(text=kandidat, images=gambar, return_tensors="pt", padding=True)
with torch.no_grad():
    out = model(**inputs)

probs = out.logits_per_image.softmax(dim=1)[0]
urut = torch.argsort(probs, descending=True)
for i in urut:
    print("%-32s %.4f" % (kandidat[int(i)], float(probs[int(i)])))
\end{lstlisting}
\end{codewindow}

\textbf{(c) Kaitan dengan Teori.}

Hubungkan hasil Anda dengan dua gagasan dari perkuliahan: (i) pilar \emph{alignment}
beserta mekanisme \emph{cross-attention} pada Pertemuan 7, dan (ii) pernyataan bahwa
penambahan modalitas \emph{tidak otomatis} meningkatkan kinerja. Apakah kasus pencarian
produk pada tugas ini termasuk kasus yang menguntungkan bagi \emph{dual-encoder}?
Sebutkan satu syarat yang membuat pilihan Anda gagal.

% =============================================================
\clearpage
\section{Kriteria Penilaian}\label{sec:evaluation}

Tugas Anda akan dinilai berdasarkan kriteria berikut:

\renewcommand{\arraystretch}{1.5}
\begin{center}\small
\begin{tabular}{|p{11.5cm}|c|}
  \hline
  \thead{\color{ctpBlue} Kriteria} & \thead{\color{ctpBlue}Bobot} \\
  \hline
  Kebenaran teknis --- kode berjalan tanpa error, \emph{InfoNCE} dan \emph{Recall@K} dihitung dengan tepat, bobot pra-latih tidak diubah. & 35\% \\
  \hline
  Kedalaman analisis --- penjelasan menghubungkan hasil dengan konsep \emph{alignment}, suhu, dan pola arsitektur multimodal. & 30\% \\
  \hline
  Eksperimen --- pengukuran baseline, perbandingan sebelum/sesudah penyelarasan, dan uji kegagalan komposisional. & 20\% \\
  \hline
  Kualitas presentasi --- tabel hasil rapi, laporan terstruktur, dan interpretasi jelas. & 15\% \\
  \hline
  \textbf{Total} & \textbf{100\%} \\
  \hline
\end{tabular}
\end{center}

% =============================================================
\section{Apa yang Dikumpulkan}\label{sec:submit}

Seluruh artefak tugas dikemas ke dalam \textbf{satu berkas arsip} dengan format penamaan:
\begin{center}
  \code{mgg09\_tugas\_nim.zip}
\end{center}

dengan \code{nim} diganti oleh Nomor Induk Mahasiswa (NIM) masing-masing tanpa spasi.

\textbf{Artefak yang dapat disertakan di dalam arsip} (dapat lebih dari satu jenis berkas, disesuaikan dengan bentuk penyelesaian Anda):
\begin{borderedsquare}
  \item \code{mgg09\_tugas\_NIM.ipynb} --- apabila penyelesaian dikerjakan menggunakan Jupyter Notebook (termasuk Google Colab). Notebook \textbf{wajib telah dijalankan} sehingga seluruh \emph{output} sel tersimpan di dalam berkas.
  \item \code{mgg09\_tugas\_NIM.py} --- apabila penyelesaian dikerjakan dalam bentuk skrip Python.
  \item \code{mgg09\_tugas\_NIM.pdf} --- apabila analisis ditulis dalam dokumen teks terpisah.
  \item \code{mgg09\_tugas\_xx.txt} --- catatan probabilitas kandidat teks pada Soal 5b; \code{xx} merupakan nomor urut (\code{01}, \code{02}, dan seterusnya).
  \item \code{mgg09\_tugas\_xx.jpg} --- apabila terdapat plot atau gambar pendukung tambahan.
\end{borderedsquare}

Setiap berkas menggunakan awalan \code{mgg09\_tugas\_NIM} yang sama, dengan \code{NIM} ditulis secara konsisten (huruf besar/kecil mengikuti NIM resmi). Pastikan seluruh berkas berada tepat di dalam satu arsip \code{.zip} tanpa struktur folder bertingkat yang tidak perlu.

\begin{callout}[ctpBlue]
\textbf{\color{ctpBlue}Catatan Pengerjaan}
\begin{itemize}
  \item \textbf{Soal 1--4:} Wajib dikerjakan oleh semua mahasiswa. Soal 1 wajib dilaporkan lebih dulu sebagai baseline.
  \item \textbf{Soal 5:} Wajib dikerjakan; bagian (b) memerlukan kode, bagian (a) dan (c) bersifat analitis.
  \item \textbf{Bobot pra-latih tidak boleh diubah:} mengubah atau melatih ulang model akan dinilai sebagai kesalahan prosedur.
  \item \textbf{Estimasi waktu:} 3--4 jam.
\end{itemize}
\end{callout}

\section*{Referensi}
{\small\color{ctpSubtext}
\begin{enumerate}[label={[\arabic*]}, leftmargin=2em]
  \item Radford, A., et al. (2021). \emph{Learning Transferable Visual Models From Natural Language Supervision}. ICML. --- sumber CLIP dan tujuan kontrastif \emph{InfoNCE}.
  \item Sch\"onemann, P. H. (1966). \emph{A Generalized Solution of the Orthogonal Procrustes Problem}. Psychometrika. --- dasar penyelarasan bentuk tertutup pada Soal 2.
  \item Vaswani, A., et al. (2017). \emph{Attention Is All You Need}. NeurIPS. --- arsitektur dasar yang mendasari sesi ini.
  \item Liang, P., Zadeh, A., \& Morency, L.-P. (2023). \emph{Tutorial on Multimodal Machine Learning}. ICML. --- bagian \emph{multimodal transformers} dan \emph{pre-training}.
\end{enumerate}}

\end{document}
